\section{Results}
\subsection{Experimental Setup}
\subsubsection{Flight-Delays Dataset}
We use a dataset derived from the popular flight-delays dataset. This dataset contains records of past flights with attributes such as origin airport, destination airport, delay reason, etc. This dataset was also used by \cite{atena2020sigmod} and we use their simulator to generate sequences as close to human analyst as possible. We use around $\sim 500k$ rows from this dataset to train the simulator. Then 6000 unique sequences are generated for this dataset by the simulator and we train our RL agent using two strategies. In the first, we train the RL agent on all generated sequences and in the second, we holdout a set of 1200 sequences and train our RL agent on remaining 4800 sequences.

\begin{figure*}[t]
\centering
\begin{minipage}{1.0\textwidth}
%\includegraphics[width=1.0\linewidth]{figures/intent_divergence.png}
\includegraphics[width=0.7\linewidth]{figures/junk_topic_model.png}
     \caption{Results of topic modeling on Flights data.}
     %\label{fig:top_distribution}
\end{minipage}
\begin{minipage}{1.0\textwidth}
%\includegraphics[width=1.0\linewidth]{figures/intent_divergence.png}
\includegraphics[width=1.0\linewidth]{figures/junk_motivation.png}
     \caption{Results showing how the distribution of topics change if EDA is used on sampled data.}
     %\label{fig:top_distribution}
\end{minipage}
\end{figure*}
\begin{figure*}[t]
\begin{minipage}{0.48\textwidth}
\centering
\includegraphics[width=1.0\textwidth]{figures/junk_reward_comparison.png}
\caption{Comparison of RL-reward for different distance metrics}
%\label{fig:framework}
\end{minipage}%
\begin{minipage}{0.48\textwidth}
\centering
\includegraphics[width=1.0\textwidth]{figures/junk_reward.png}
\caption{Characteristics of reward for RL training.}
%\label{fig:framework}
\end{minipage}
\end{figure*}

